\begin{table}[t]
\centering
\small
\caption{Open-source candidate-provider sensitivity on the crossed/forward cohort. The candidate route refits the source models on seven candidate providers with recorded reuse paths while retaining the target masks, daily features, and plug-in baselines. Retention and coverage are D2 quantities.}
\label{tab:S6}
\resizebox{\linewidth}{!}{%
\begin{tabular}{lrr}
\toprule
Metric & Original source & Candidate source \\
\midrule
D2 trend RMSE ($^{\circ}$C/10a) & 0.164 & 0.158 \\
Classical trend RMSE ($^{\circ}$C/10a) & 0.147 & 0.156 \\
D2 retention in the symmetric MSE decomposition (\%) & 93.9 & 93.7 \\
Nominal 95\% D2 interval coverage (\%) & 88.0 & 87.6 \\
Common four-way pairs & 738 & 738 \\
Connected groups & 17 & 17 \\
Factor cases for retention & 807 & 807 \\
Coverage cases & 820 & 820 \\
\bottomrule
\end{tabular}}
\par\smallskip\footnotesize Notes: The candidate route leaves the ranking of D2 and the classical comparator unchanged. It is a candidate open reproducibility route, not a legal clearance: LUBW rights and the provenance of neighbor records still require confirmation. Candidate-source tables are descriptive, repetitions are not independent climates, and the filtered inputs remain processed study data rather than provider-native records. Source: \path{results/ems_revision_20261102/flow_diagnostic.json}, extracted from \path{paper/wrr_tr/open_source_values.tex}.
\end{table}
